\section{Clearance-to-roof coarea without a fixed-clearance density bound}
\label{sec:cr-coarea}
We keep the physical parameters and every event in Section~1. In this
section $u$ is the clearance coordinate of Section~\ref{sec:cf-flux},
not physical time. Physical time is denoted by $t$ throughout.
Fix $R,m,n,k$, a flight $0\le j<m$, a sign $\sigma$, and an actual
regular marked itinerary. On a local angular lift put
\[
 \mathcal Q_j(\vartheta,u)
   =(\alpha_0,\alpha_m)(X_{R,j}(\vartheta,\sigma,u)),
 \qquad S(\vartheta,u)=F(\mathcal Q_j(\vartheta,u)).
\]
The point used as a mark is an unvisited near miss, not an additional
reflection. Every derivative below is taken on the same designated
regular itinerary. In particular $m$, $j$ and the physical roof are
not changed by this parametrization.

\begin{theorem}[The long-record Jacobian and its critical Hessian]
\label{thm:cr-jacobian}
Let $A=D\mathcal Q_j$, let $g=\nabla_q F$, and let $G=D_q^2F$.
On every such chart,
\begin{align}
 |\det A|&=(R\omega)^{-1},\label{eq:cr-area}\\
 |\det D_{(\vartheta,u)}(S,\arg g)|
    &=\chi/(R\omega) \quad (g\ne0),\label{eq:cr-gauss}\\
 D^2S&=A^{\mathsf T}GA+
                \sum_{a=1}^2g_aD^2(\mathcal Q_j)_a.
                  \label{eq:cr-chain-hessian}
\end{align}
At a regular critical roof, $g=0$, the Hessian $D^2S$ is positive
definite and
\begin{equation}\label{eq:cr-critical-determinant}
 \sqrt{\det D^2S}=\frac{\sqrt{\det G}}{R\omega}.
\end{equation}
These are identities for each finite physical collision count, with
all coordinate and section-normalization factors displayed. In
particular the Hessian term containing $g$ cannot be dropped away
from a critical roof.
\end{theorem}
\begin{proof}
The collision map preserves $d\alpha\,dp$ at fixed $R$. The signed
line map at collision $j$ has absolute Jacobian $R$, by
\eqref{eq:cf-jacobian}. Composition with $T_R^j$ therefore still has
absolute Jacobian $R$ with respect to the initial collision
coordinates. The map $(\alpha_0,p_0)\mapsto(\alpha_0,\alpha_m)$ has
absolute Jacobian $\omega^{-1}$ by the definition of $\omega$.
Dividing the two determinants gives \eqref{eq:cr-area}. It uses the
actual inverse iterate and is independent of where the mark occurs.
On a noncritical chart the determinant of
$q\mapsto(F,\arg g)$ is $\chi$; the chain rule gives
\eqref{eq:cr-gauss}. Twice differentiating $F\circ\mathcal Q_j$
gives \eqref{eq:cr-chain-hessian}, component by component. At $g=0$
its last term vanishes, and the retained positive-definite optical
Hessian and the invertibility of $A$ give the assertion. Taking the
determinant and using \eqref{eq:cr-area} proves
\eqref{eq:cr-critical-determinant}.

As a normalization check, the line-coordinate density is
$(4\pi Rc_*)^{-1}$. Dividing it by the determinant in
\eqref{eq:cr-gauss} gives exactly $\omega/(4\pi c_*\chi)$, the
original target-roof/Gauss-angle density. There is no additional
factor $R$, $R^{-1}$ or $c_*^{-1}$. All chart restrictions and exact
labels remain on the source side of this change of variables.
\end{proof}

\begin{corollary}[Uniform determinant growth on regular marked charts]
\label{cor:cr-determinant-growth}
With the retained constants $\varrho=47/59$ and $\lambda_\circ$,
all regular marked charts with the original endpoint cap satisfy
\[
 |\det D\mathcal Q_j|\ge\frac{600}{2209}\varrho^{-(m-1)},
 \qquad
 |\det D(S,\arg g)|\ge
       \frac{600\lambda_\circ}{2209}\varrho^{-(m-1)}
       \quad(g\ne0).
\]
At a critical roof the same lower bound as in the second inequality
holds for $\sqrt{\det D^2S}$. The constants are uniform in the
original radius interval, the exact labels, the flight index, and
the physical word. These determinant bounds do not bound individual
singular values, a chart's condition number, or the complete
physical sum of inverse determinants.
\end{corollary}
\begin{proof}
The retained optical estimate is
$\omega\le(47/6)\varrho^{m-1}$, and $R\le47/100$.
Insert these in \eqref{eq:cr-area}. Then use
$\chi\ge\lambda_\circ$ and
$\sqrt{\det G}\ge\lambda_\circ$ in
\eqref{eq:cr-gauss} and \eqref{eq:cr-critical-determinant}.
The identity $6/(47(47/100))=600/2209$ fixes the common constant.
\end{proof}

The following elementary estimate is recorded with proof to avoid
importing an unstated finite-type or regularity theorem.

\begin{lemma}[A one-dimensional band estimate]
\label{lem:cr-band}
Let $f\in C^\ell(I)$ on an interval, with integer $\ell\ge1$ and
$|f^{(\ell)}|\ge\kappa>0$. For $h>0$,
\begin{equation}\label{eq:cr-band}
 |\{x\in I:a<f(x)<a+h\}|
          \le 2\ell(h/\kappa)^{1/\ell}.
\end{equation}
The estimate also bounds any Borel subset of that band.
\end{lemma}
\begin{proof}
First restrict the band to a compact subinterval and let its measure
be $L>0$. Its cumulative Lebesgue measure is continuous and
one-Lipschitz. The $\ell+1$ quantiles with successive mass gaps
$L/\ell$ can be chosen in its closed support; their successive
coordinate gaps are at least $L/\ell$. Equivalently one can use
interior approximate quantiles and pass to their limits. Write them
as $x_0<\cdots<x_\ell$. Continuity puts $f(x_i)$ in $[a,a+h]$.
The divided-difference mean-value theorem, proved by subtracting the
interpolating polynomial and applying Rolle's theorem $\ell$ times,
gives some $\xi\in(x_0,x_\ell)$ for which
\[
 \frac{f^{(\ell)}(\xi)}{\ell!}
  =\sum_{i=0}^{\ell}
       \frac{f(x_i)-a}{\prod_{v\ne i}(x_i-x_v)}.
\]
The subtraction of $a$ is legitimate because the divided difference
of a constant is zero. The gap bound makes the absolute denominator
at least $(L/\ell)^\ell i!(\ell-i)!$. Hence
\[
 \kappa\le h(\ell/L)^\ell
       \sum_{i=0}^{\ell}{\ell\choose i}
       =h(2\ell/L)^\ell.
\]
Rearrange, then exhaust the original interval by compact
subintervals. This also handles a band not initially known to have
finite measure. No statement about the number of physical words is
used in the proof.
\end{proof}

\begin{theorem}[A positive finite-type clearance slice]
\label{thm:cr-finite-type}
Suppose $S\in C^{\ell+1}$ on a neighbourhood of
$I\times[0,\delta_0]$, and
\[
 0<b_f\le |S_u|\le B_f,\qquad
 |\partial_\vartheta^\ell S(\vartheta,0)|\ge\kappa_f>0.
\]
The sign of $S_u$ is fixed on this rectangle. Let $E$ be any Borel
subset of it and $0\le h\le H_f$ a pointwise bounded Borel weight.
For $0<\delta\le\delta_0$, the measure
\[
 S_\#\bigl(h\one_E\one_{\{0<u<\delta\}}
                         \,d\vartheta\,du\bigr)
\]
has a density $D_{E,h,\delta}$ with
\begin{equation}\label{eq:cr-finite-type}
 \|D_{E,h,\delta}\|_\infty
 \le\frac{2\ell H_f}{b_f}
                  \left(\frac{B_f\delta}{\kappa_f}\right)^{1/\ell}.
\end{equation}
In the physical marked record one takes
$h=(4\pi Rc_*)^{-1}$, or a bounded multiplier of this weight, and
$E$ is the original physical domain with all exact labels. The
constants must be established on that actual chart; the theorem
does not assert their uniformity in $m$ or their summability over
all charts.
\end{theorem}
\begin{proof}
At fixed $\vartheta$, $u\mapsto S(\vartheta,u)$ is strictly
monotone. Let $U_t(\vartheta)$ denote its unique inverse when it
exists. The ordinary one-dimensional substitution, first against
nonnegative roof tests and then by Tonelli, gives
\[
 D_{E,h,\delta}(t)=\int_I
 \frac{h(\vartheta,U_t)\one_E(\vartheta,U_t)
                 \one_{\{0<U_t<\delta\}}}
      {|S_u(\vartheta,U_t)|}\,d\vartheta.
\]
Missing inverses contribute zero. This argument proves absolute
continuity, not merely an estimate conditional on its existence.
If $S_u>0$ and the inverse lies in $(0,\delta)$, then
$t-B_f\delta<S(\vartheta,0)<t$. For $S_u<0$ the corresponding
band is $t<S(\vartheta,0)<t+B_f\delta$. Both have width
$B_f\delta$. Bound the integrand by $H_f/b_f$ and apply
Lemma~\ref{lem:cr-band}. Borel cuts, including section cuts, only
restrict the integral and never create a chart-count multiplier.
The same proof on bounded compact rectangles followed by monotone
exhaustion gives the assertion for open intervals and Borel weights.
\end{proof}

\begin{theorem}[Integrated simple folds, with a sharp local constant]
\label{thm:cr-fold}
Assume that on $I\times(0,\delta_0)$ one has
$S_{\vartheta\vartheta}\ge a_f>0$ and $S_u\ge b_f>0$, and that
for each $u$ the unique angular critical point
$\vartheta_0(u)$ exists in $I$. For the same positive weights and
selectors as above,
\begin{equation}\label{eq:cr-fold}
 \|D_{E,h,\delta}\|_\infty
                 \le 2H_f\sqrt{\frac{2\delta}{a_fb_f}}.
\end{equation}
The assertion also applies after reversing the roof sign or the
clearance-coordinate orientation when the corresponding derivative
signs are reversed. It does not require a bounded roof density at
each fixed clearance.
\end{theorem}
\begin{proof}
Put $v(u)=S(\vartheta_0(u),u)$. The implicit function theorem and
the critical-point equation give $v'(u)=S_u(\vartheta_0(u),u)
\ge b_f$. For fixed $u$ there are at most two angular roots of
$S(\vartheta,u)=t$, and
\[
 |S_\vartheta(\vartheta,u)|^2
                \ge2a_f(S(\vartheta,u)-v(u)).
\]
Indeed differentiate $S_\vartheta^2$ along each monotone side and
integrate from its zero at $\vartheta_0$. On the left side the
integration orientation and the sign of $S_\vartheta$ reverse
together, giving the same inequality. One-dimensional substitution
on those two sides gives a fixed-$u$ density bounded by
$H_f\sqrt{2/a_f}(t-v(u))^{-1/2}$ for $t>v(u)$ and zero below.
The angular critical point itself has zero angular measure. The
bound is allowed to be infinite at $t=v(u)$.

For $0<u<\delta$, if $t$ is in the range of $v$, let $u_t$ be its
unique inverse. The contributing values satisfy $u<u_t$, and
$t-v(u)\ge b_f(u_t-u)$. If $t\ge v(\delta)$, replace $u_t$
by $\delta$ for an upper bound; if $t\le v(0)$ there is no
contribution. In every case
\[
 \int_0^\delta\one_{\{t>v(u)\}}
                     (t-v(u))^{-1/2}\,du
                         \le2\sqrt{\delta/b_f}.
\]
Integrate the fixed-$u$ estimate by Tonelli to obtain
\eqref{eq:cr-fold}. Restricting by $E$ preserves the inequality.
For $S(\vartheta,u)=\vartheta^2+u$, $h=1$, and the whole real
angular line, the integrated density is
\[
 2\{\sqrt{t_+}-\sqrt{(t-\delta)_+}\},
\]
whose maximum is $2\sqrt\delta$. Thus the constant is attained
when $a_f=2$ and $b_f=1$. This elementary normal form checks the
local analytic bound only; it is not substituted for a billiard
trajectory or for the joint-record distribution.
\end{proof}

\subsection{One-copy summation and the remaining uniformity}
On a countable collection of admissible physical charts, first
replace their source sets by disjoint Borel differences. The
preceding density formulae then add exactly. Their positive upper
bounds add as well, with an extended value permitted until a
finite bound has been proved. Two impact signs and distinct
unvisited centres are retained with precisely the multiplicity in
\eqref{eq:cf-record-disintegration}; no full-word cardinality is
inserted. Borel event cuts do not require a new chart merely because
they disconnect the selected set.

In particular the per-clearance condition
\eqref{eq:cf-bridge-input} is not needed on a chart controlled by
Theorem~\ref{thm:cr-fold}. An angular fold may have an unbounded
conditional density and a bounded integrated density at the same
time. What is still required for the full clearance payment is a
bound for the complete, exact-label, positively weighted chart sum,
in the original $m^2 L^\infty(dt)$ topology. The local constants in
\eqref{eq:cr-finite-type} or \eqref{eq:cr-fold} have not been
proved to provide that bound. Charts with $S_u=0$, degenerate
angular behaviour outside a verified finite type, and singular
physical boundaries are not discarded. Regular full critical roofs
are treated next by a different exact calculation.
